% Compatibility entry point; active section source: concepts.tex.
\section{Key concepts}
\label{sec:concepts}

\subsection{Image and caption retrieval}
An \emph{embedding} is a list of numbers that represents an image or a text caption. Embeddings from a shared model can be compared to find items with similar representations. In this project, the image encoder is used to retrieve training images similar to a held-out query image. The caption index stores vectors for training captions, but a caption-generation pipeline that uses these results has not been implemented.

\subsection{What similarity does and does not mean}
The project uses CLIP, a vision-language model trained to align image and text representations \citep{radford2021clip,clipmodelcard}. A high similarity score means that two vectors are close under the chosen representation. It does not prove that the retrieved image is an exact match, that its caption applies to the query, or that a future generated caption is factually correct.

FAISS is a library for efficient similarity search over vectors \citep{douze2024faiss,faissdocs}. The project uses its exact flat inner-product index. Because the stored vectors are L2-normalized, inner product gives cosine similarity. This index provides a straightforward reference search, not an approximate search optimization.

\subsection{Phase boundaries}
Phase 1 ends with saved splits, analysis, embeddings, indexes, and retrieval checks. Phase 2 must add a generator, define retrieval conditions and prompts, run controlled generation, and evaluate the outputs. The two phases should not be confused when interpreting the results in this report.

\endinput
\chapter{Background and Assigned Technique}\label{ch:sota}
\section{Retrieval-augmented generation}
RAG combines a retriever over non-parametric memory with a generator that conditions on retrieved material. In the original formulation, dense retrieval selects passages and a sequence-to-sequence model conditions on them for knowledge-intensive NLP tasks \cite{lewis2020rag}. Subsequent work has explored retrieval over examples, documents, demonstrations, and multimodal content. For image captioning, retrieved items may be similar images, their captions, structured annotations, or combinations of these. These choices are not interchangeable: image-neighbor retrieval can provide visual analogies, while caption retrieval can offer lexical templates but may also transfer incorrect objects or relations.

The project implements vector indexes that could support example retrieval, but it does not yet implement a generator that consumes retrieved records. Consequently, the current system is a retrieval-capable dataset pipeline rather than a complete RAG model. A meaningful comparison must specify retrieval query, candidate corpus, top-$k$, filtering, prompt construction, generator, decoding, and controls.

\section{Image captioning and evaluation}
COCO has been widely used to study object recognition and image captioning in complex everyday scenes \cite{lin2014coco}. Its multiple captions per image capture some of the diversity of human descriptions, but references are not exhaustive descriptions of every visible fact. N-gram metrics such as BLEU and ROUGE-L measure overlap with references; METEOR adds stemming and synonym-aware matching; CIDEr emphasizes consensus with multiple references; SPICE compares semantic scene-graph propositions; learned similarity measures such as BERTScore compare contextual token representations \cite{papineni2002bleu,banerjee2005meteor,lin2004rouge,vedantam2015cider,anderson2016spice,zhang2020bertscore}. Each metric captures a different aspect and none alone establishes visual truth. The repository has not generated captions and reports none of these scores.

\section{Vision-language representations}
CLIP learns aligned image and text spaces through contrastive image-text pretraining and supports transfer using natural-language supervision \cite{radford2021clip}. This makes CLIP embeddings a practical common representation for image and caption similarity. Similarity is not equivalent to caption correctness: semantically related images may differ in small objects, counts, actions, or spatial relations. In this project CLIP is used for embedding and zero-shot exploratory predictions; those predictions are not ground-truth annotations.

\section{Assigned Technique and Implemented Method}\label{ch:method}
The assignment's exact designated RAG technique cannot be identified from the source PDF because that file was not found. The supplied task text names a comparative RAG study but does not provide a verifiable technique assignment. Rather than invent an assigned method, this report describes the repository's actual method: dense cosine retrieval using CLIP embeddings and exact inner-product search, with separate image and caption memories. No caption-generation strategy has been selected or implemented.

For an image representation $q$ and training representation $v_i$, retrieval ranks candidates by cosine similarity
\begin{equation}\label{eq:cosine}
 s(q,v_i)=\frac{q^\top v_i}{\lVert q\rVert_2\lVert v_i\rVert_2}.
\end{equation}
Embeddings are L2-normalized, so $s(q,v_i)=q^\top v_i$. The indexes use FAISS \texttt{IndexFlatIP}, a flat exact inner-product index. For a future comparison, this dense-neighbor baseline can be contrasted with caption-neighbor retrieval, hybrid image-plus-caption retrieval, and a no-retrieval generator, provided each method uses the same generator and held-out examples.
